% script.tex -- rehearsal speaker script for talk/talk.tex (40-minute total session).
% Contract: talk/structure-plan.md (sections 2-4, 8, 10, 11) and talk/notes/audit-adoptions.md.
% One block per main frame, keyed to the exact frame title in talk.tex; the duplicated build
% (F10a -> F10b) has one block per physical frame and an \Advance mark at the end of the first.
% Spoken numbers are rounded forms of the numbers on the slides and carry the same registry key
% in a "% source:" comment (numerics/quantity_registry.csv); none is stronger than its slide.
% Check and compile:
%   python check_script.py script.tex --deck talk.tex --slot-minutes 40
%   python compile_deck.py script.tex --themes-dir talk
\documentclass[11pt]{article}
\usepackage[margin=0.78in,headheight=14pt]{geometry}
\usepackage{xcolor}
\usepackage{microtype}
\usepackage{needspace}
\usepackage{setspace}
\usepackage{fancyhdr}
\usepackage{amsmath,amssymb}

\newcommand{\PaperTitle}{Competition Creates Competition: Stock Prices and the Discovery of Takeover Bidders}
\newcommand{\PaperShortTitle}{Competition Creates Competition}

\definecolor{ScriptBlue}{HTML}{1F4E79}
\definecolor{ScriptGray}{HTML}{666666}
\setstretch{1.12}
\setlength{\parindent}{0pt}
\setlength{\parskip}{0.45em}
\pagestyle{fancy}
\fancyhf{}
\lhead{\small\textcolor{ScriptGray}{\PaperShortTitle}}
\rhead{\small\textcolor{ScriptGray}{Rehearsal script, 40-minute session}}
\cfoot{\thepage}

\newcommand{\Pause}{\textcolor{ScriptGray}{\textit{[pause]}}}
\newcommand{\Advance}{\textcolor{ScriptGray}{\textit{[advance]}}}
\newcommand{\Cue}[1]{\par\nobreak\smallskip\noindent
  \textcolor{ScriptGray}{\footnotesize\textbf{Live cue:} #1}\par}
\newcommand{\scriptopening}[2]{%
  \Needspace{10\baselineskip}%
  \section*{Opening \hfill\normalfont\small\textcolor{ScriptGray}{#1}}
  #2\par}
\newcommand{\scriptframe}[3]{%
  \Needspace{9\baselineskip}%
  \par\medskip\noindent\textcolor{ScriptBlue}{\rule{\linewidth}{0.55pt}}%
  \par\smallskip\noindent
  \textbf{#1}\hfill\textcolor{ScriptGray}{\small #2}\par\smallskip
  #3\par}
\newcommand{\ScriptAppendix}{%
  \clearpage
  \section*{\mbox{Q\&A appendix} \hfill
    \normalfont\small\textcolor{ScriptGray}{conditional; excluded from talk time}}}

\title{Speaker script\\[0.25em]\large \PaperTitle}
\author{Austin Li}
\date{40-minute session: about 30 minutes of prepared speech, about 10 for questions}

\begin{document}
\maketitle
\thispagestyle{fancy}

{\small\textcolor{ScriptGray}{Spoken symbols: ``r-zero, r-one, r-two'' for the weak, strong and
stronger-still incumbent; ``Delta-T'' for the target-payoff spread; ``c-L, c-H'' for the cheap and
expensive preparation cost. Checkpoints (structure plan, section 10): frame 3 by 4.75 min; frame 5 by
9.5 min, else skip frame 6 and say its two numbers on frame 5; frame 9 by 18.5 min, else skip frames
12 and 14; frame 11 by 23.5 min, else skip frame 15 and give frame 13 as one sentence.}\par}

% ======================================================================
\scriptopening{about 0.4 min}{%
  Thank you for having me. I'm Austin Li, and the paper is \emph{\PaperTitle}.
  It is a single-authored theory paper about who shows up to compete for a
  takeover target, and what the target's stock price has to do with it. Please
  ask questions as we go.
  \Cue{Thank you for having me; single-authored theory paper; questions welcome.}
}

% ======================================================================
\scriptframe{Does a stronger incumbent keep the challenger out?}{about 2.3 min}{%
  Here is the situation the paper is about. Read the three boxes from left to
  right. An approach or a strategic review has been disclosed, so the target is
  publicly in play. Its stock keeps trading. And a second bidder is deciding
  whether to get ready to bid.

  One bidder, the \emph{incumbent}, is already prepared. The market knows how
  strong it is, in the sense of the distribution of its value, not the bid it
  will make. The second bidder, the challenger, cannot just show up with a
  number. To make an executable bid it has to pay for diligence, financing and
  approvals. I will call paying that cost entry. A challenger that declines
  simply stays out.

  So does a stronger incumbent keep the challenger out? The received answer is
  yes. Fishman, and Hirshleifer and Png, give the logic: a stronger rival raises
  what a winner has to pay, so preparation earns less, and fewer challengers
  prepare. I am not going to tell you that logic is wrong. It holds whenever
  what the challenger knows is fixed.

  But here the challenger decides while the stock trades, and it can watch the
  price. The model needs an interval, not a date. A disclosed approach or an
  open contest can create one; a wholly confidential process does not. The
  Imprivata proxy shows the stages separately: an unsolicited approach,
  outreach to other potential buyers, and indications of interest conditional
  on more diligence. That is a costly preparation stage in the record. It is
  not evidence that a price drew anyone in.

  The twist is on the last line. That price is an equilibrium object, and what
  it reveals depends on how strong the incumbent is. So my claim is not that
  deterrence is wrong. It is that deterrence can be \emph{outweighed}, and I owe
  you the set of primitives on which that happens.
  \Cue{Here is the situation the paper is about.}
}

% ======================================================================
\scriptframe{This paper}{about 2.2 min}{%
  Here is what I do and what I find. I study a takeover auction in which an
  informed investor trades the target's stock before the challenger decides
  whether to prepare. Prices are rational: market makers price the entry that
  each price will induce. Trading and entry are both solved in equilibrium.

  The main result is the title. Competition creates competition. A stronger
  incumbent can turn an uninformative stock price into an informative one, and
  that raises entry. In the benchmark economy entry roughly doubles, from a
  quarter to just over one half.
  % source: base_entry_weak (0.250), base_entry_strong (0.523)
  And that happens while the challenger's expected profit before any news
  falls, from about four point eight to about four point three.
  % source: base_profit_prior_weak (4.80), base_profit_prior_strong (4.29)
  \Pause{} The challenger keeps \emph{less}, and more challengers come.

  Why? One auction, two claims. The sale rule splits one surplus. The
  challenger holds what it keeps if it wins; target shareholders hold what the
  auction pays them. A strong incumbent can beat a low-value challenger but not
  a high-value one. So when the incumbent gets stronger, the challenger keeps
  less, but the sale price depends more on who the challenger is. That makes the
  investor's information worth trading on. Trading puts it into the price, and
  good news in the price is what justifies expensive preparation.

  The second claim: information is the channel. Hold fixed what the price
  reveals, and a stronger incumbent cannot raise entry. That is the textbook
  result, and it is a theorem in this model too.

  The implication is on the bottom line. When the target trades, a stronger
  rival is not a pure deterrent, because who competes depends on what the
  price reveals before anyone prepares.

  Here is the route: the model, the two forces, the theorem, the numbers, and
  then the control that shows information does the work.
  \Cue{Here is what I do and what I find.}
}

% ======================================================================
\scriptframe{What is new relative to learning from prices}{about 0.7 min}{%
  Three literatures meet here. From the feedback literature, Dow, Goldstein and
  Guembel, and Edmans, Goldstein and Jiang, I take the idea that prices guide
  real decisions. What the sale rule adds is a split. One surplus divides into
  a traded claim and the challenger's claim, and competition moves them in
  opposite directions. From Fishman I keep deterrence intact; I add a market
  that trades before preparation. From Levin and Smith, entry into auctions is
  endogenous; here it responds to what the price reveals. The increment is that
  opposition between the two claims, and the entry reversal it produces.
  \Cue{Three literatures meet here.}
}

% ======================================================================
\scriptframe{Model: who moves, who knows what, and when}{about 2.8 min}{%
  Here is the model, in one timeline. Read the order of moves once, because it
  is the load-bearing assumption.

  Box one. The seller commits to a cash second-price auction with reserve $p$.
  Commitment comes before any trading, so the seller is not reacting to the
  price.

  Box two. An investor knows the challenger's quality, theta, and trades the
  target's stock against noise traders. The investor has no position, no
  control rights and no way to bid. It knows theta, not the incumbent's value.
  Neither bidder trades.

  Box three. Market makers see only total order flow. They price the share at
  expected sale proceeds, and that includes the entry each price will induce.
  So nothing here rests on mispricing.

  Box four. The challenger sees the price and its own preparation cost. It never
  sees the order flow, theta, or the incumbent's value. Entry means paying the
  cost; that is what reveals theta and lets it bid. Without preparation there is
  no executable bid, so a challenger that declines stays out.

  Box five. Both bidders bid truthfully. The incumbent's value, capital R, is
  uniform between zero and small r, and I will call small r incumbent strength.
  It is public. It is a distribution, not a bid. The challenger is worth h or
  ell, equally likely.

  Now the first legend line. The ordering is reserve $p$, below ell, below
  strength, below h. So the incumbent can beat a low-value challenger but never
  a high-value one. Hold on to that; the next slide lives on it. And $p$ is the
  reserve, not the stock price.

  The green line is the preparation cost. It is cheap, c-L, with probability
  rho, and expensive, c-H, otherwise, independent of the challenger's value. The
  two costs have different jobs. Cheap preparation gives an entry floor.
  Expensive preparation is the margin a good price can recruit.

  If preparation came before trading, there would be nothing to learn. That
  order is the whole point.

  You may be wondering why an investor would know theta when the challenger
  does not. The benchmark gives the investor perfect information only to make
  the mechanism easy to see. Section 5.3 removes both extremes, and the result
  survives a challenger whose own signal is more accurate than the investor's.
  \Cue{Here is the model, in one timeline.}
}

% ======================================================================
\scriptframe{One auction, two claims, opposite responses}{about 2.3 min}{%
  This table is the whole paper in one place. It is the cash second-price
  auction, once the challenger has prepared. Read the rows by the incumbent's
  realized value, capital R.

  Top row: R below ell. Either type of challenger wins and pays the larger of
  the reserve and R. Target shareholders get the same amount whoever the
  challenger is, so the gap in the last column is zero.

  Bottom row: R above ell. Now the types split. A high-value challenger still
  wins, but pays R. A low-value challenger loses; the incumbent wins and pays
  ell, the low-value challenger's bid. So a \emph{losing} low-value challenger
  sets what the incumbent pays, and the gap in target proceeds is R minus ell.
  That asymmetry is the source of everything that follows.

  Now make the incumbent stronger, which shifts R up. Two things happen.
  \Pause{} First, the orange line. The target-payoff spread, Delta-T, is the
  difference in expected target proceeds between a high-value and a low-value
  challenger. It is the expected gap from the last column, so it rises: more of
  the incumbent's values sit above ell. Second, the black line. The challenger's
  gross profit falls, because what it pays when it wins goes up. That holds for
  both types.

  So one shift has opposite signs. The challenger keeps less. Target proceeds
  depend more on who the challenger is. Proposition 1 says both signs hold for
  any first-order strengthening of the incumbent, not just the uniform case.

  You might ask: why a second-price auction? With private values, an open
  ascending contest ends at the runner-up's value, just as this rule does, which
  is why takeover-contest models commonly use it. The property that matters is
  that a losing low-value challenger sets the incumbent's price. It is not free.
  Under bargaining the spread effect can flip sign, and there is a backup on
  that. First-price equilibria I do not solve.
  \Cue{This table is the whole paper in one place.}
}

% ======================================================================
\scriptframe{A stronger incumbent: spread up, profit down}{about 1.6 min}{%
  Here are the two claims at benchmark scale. Both panels put incumbent
  strength on the horizontal axis. The dotted lines mark the two economies I
  will compare: a weak incumbent, r-zero, at 1.2, and a strong one, r-one, at 3.
  % source: base_r_weak (1.2), base_r_strong (3)

  Panel (a), in orange, is the target-payoff spread. Against the weak incumbent
  it is almost nothing, because the incumbent rarely draws a value above ell.
  It climbs steadily with strength, to about two thirds at r-one.
  % source: base_spread_weak (0.0167), base_spread_strong (0.667)

  Panel (b) is the challenger's expected profit at the prior, before any news.
  It slopes down, gently, from about four point eight to about four point three.
  % source: base_profit_prior_weak (4.80), base_profit_prior_strong (4.29)

  Notice the proportions. The spread moves by a large multiple; profit moves by
  about a tenth. But the size of the profit fall is not what the result turns
  on. What matters is whether the spread clears the cost of trading on it.

  These are auction-stage payoffs only. Nothing here is an equilibrium of
  trading or entry yet.

  So deterrence holds at every fixed belief: a challenger with no news wants to
  prepare less against the strong incumbent. But there is now much more for the
  stock to reveal. The spread goes from almost nothing to something worth
  trading on. The next question is whether anyone trades on it.
  \Cue{Here are the two claims at benchmark scale.}
}

% ======================================================================
\scriptframe{Trading pays only against a strong incumbent}{about 3.0 min}{%
  Now the investor. The question is when informed trading pays. After seeing
  theta, the investor places an order: q-H after a high value, q-L after a low
  value. Orders are bounded, each unit costs k to trade, and prices anticipate
  entry.

  First, what the price can say. It reveals the market's belief, mu, that the
  challenger is high value. Noise keeps that belief inside a band, from small m
  to capital M. With Laplace noise of scale two, small m is about a quarter,
  % source: base_m (0.27), base_b (2)
  so even the best order flow never makes the market certain.

  Next, what the investor earns. Market makers already price the entry each
  price will induce. What the investor keeps is the residual advantage: what it
  knows a share pays, minus its price. The table breaks that into three
  factors. The probability the challenger enters, times the target-payoff
  spread, times the market's remaining uncertainty. \Pause{} Entry is at least
  rho, because cheap preparation pays after any price; that is the low-cost
  floor on the next slide. Remaining uncertainty is at least m. So the advantage
  sits between rho times m times Delta-T, and Delta-T itself.

  Now the two economies. Against the weak incumbent the advantage is at most the
  spread, which is under two hundredths. The trading cost is two hundredths.
  % source: base_spread_weak (0.0167), base_k (0.02)
  Every order loses money, so the unique outcome is no trade.

  Against the strong incumbent, each extra unit earns at least that lower bound,
  times a haircut of one half for the order's own price impact. With Laplace
  noise, one more unit erodes the per-unit advantage by at most a fraction one
  over b. That bound comes out just above k.
  % source: base_k + base_margin_strong_trade (0.0224 > 0.02)
  So every marginal unit pays, and the investor trades fully: buy one unit after
  good news, sell one after bad.

  That is the trading-cost window in the gray line: k sits above the weak spread
  and below the strong bound. The margin on the strong side is small at the
  benchmark; I will say what that does and does not mean when we get to the
  theorem.

  Two things this is not. It is not a toehold story: the investor has no
  position, no control rights and no route to bidding. And it is not
  manipulation. A wrong-signed order, such as buying on bad news to fake good
  news, loses against every price schedule the market could use.

  So the same shift that lowers the challenger's profit switches informed
  trading \emph{on}.
  \Cue{Now the investor.}
}

% ======================================================================
\scriptframe{Only good news makes expensive preparation pay}{about 2.3 min}{%
  Now the challenger's side: what does a moving price tell it? The figure plots
  the challenger's expected gross profit against incumbent strength, at three
  beliefs. The solid curve is the best belief a price can induce, capital M. The
  dashed curve is the prior, one half. The dotted curve is the worst belief,
  small m. The two green lines are the costs: cheap at one, expensive at six.
  % source: base_c_low (1), base_c_high (6)

  All three curves slope down. That is deterrence again, now at every belief.

  First, the low-cost floor. The dashed green line sits below even the
  worst-belief curve at r-one. So a cheap challenger enters after any price,
  which gives entry of at least rho, the bound we used on the last slide.

  Second, the high-cost window. Look at the expensive-cost line at six. Against
  the weak incumbent, profit at the prior is about four point eight, below six,
  % source: base_profit_prior_weak (4.80), base_c_high (6)
  so with no news the expensive challenger stays out. Against the strong
  incumbent, the best-price curve is above six. So good news, and only good
  news, brings the expensive challenger in.

  Think of the price as a coarse reading of the challenger's quality. The cheap
  challenger does not need it; it prepares anyway. The expensive challenger
  needs it, and only a good enough reading is worth paying six for. So the price
  acts on exactly one margin, the expensive challenger.

  Further right, the best-price curve crosses the expensive-cost line. Beyond
  that ceiling strength, even the best price falls short. The r-two marker sits
  there; hold it for later.

  Now notice something. Against the weak incumbent the best-price curve also
  clears six. That does not help. We just saw that the weak incumbent's price
  never moves, because nobody trades. What the weak incumbent lacks is not a
  good curve. It is a price that \emph{moves}. Hold that thought for the control.
  \Cue{Now the challenger's side: what does a moving price tell it?}
}

% ======================================================================
\scriptframe{Proposition 2: a stronger incumbent can raise entry}{about 2.6 min}{%
  Here is the main theorem, in words first. Take three economies that differ
  only in incumbent strength.

  Part one. Against the weak incumbent, r-zero, the unique outcome is no trade.
  The price is uninformative, and only cheap challengers enter, so entry equals
  rho.

  Part two. Against the strong incumbent, r-one, the unique outcome is full
  orders. The price is informative. Entry is strictly above rho, and a
  high-value challenger acquires the target more often. Ownership rises too
  because good news is likelier when the challenger really is high value, so
  the extra entry tilts toward high values.

  Part three. Stronger still, at r-two, the floor and profitable trading still
  hold. Only the best price stops covering expensive preparation. Trading stays
  fully informative, but entry falls back to rho.

  Why is it true? Against the weak incumbent there is nothing worth trading on,
  so the price is flat and the challenger learns nothing. Against the strong
  incumbent the spread is large enough that the investor must trade. The price
  moves, and good news pushes the belief past the point where expensive
  preparation pays. Deterrence is still there. It is just smaller than what the
  price now tells the challenger.

  The three conditions are the ones we have built, now by name. The low-cost
  floor: cheap preparation pays even at the worst belief. The high-cost window:
  expensive preparation fails at the weak prior but pays at the best strong
  price. The trading-cost window: k sits between the weak spread and the strong
  bound.

  On scope, in the gray note. Parts one and two hold on a nonempty open set of
  primitives, for every h above ell. Part three I show at the benchmark and
  nearby. Uniqueness allows arbitrary mixed orders and every unilateral
  deviation in the order interval, so no candidate profile is assumed.

  The bottom line is the chain. A stronger incumbent widens the spread; informed
  trading pays; the price becomes informative; good news clears the expensive
  cost; and the expensive challenger enters. That chain has to beat deterrence,
  which never goes away. Under these three conditions, it \emph{does}.
  \Cue{Here is the main theorem, in words first.}
}

% ======================================================================
\scriptframe{At the benchmark, entry rises from 0.250 to 0.523}{about 1.4 min}{%
  Now the numbers. Let me orient the table first. Two columns, weak and strong.
  The top two rows are the auction payoffs. The middle rows are trading and the
  price. The bottom rows are the outcomes. Entry is the probability that the
  challenger prepares; high-value ownership is the probability that a
  high-value challenger ends up owning the target.

  Start at the top. The challenger's profit at the prior falls, from about four
  point eight to about four point three. The spread rises from almost nothing
  to about two thirds.
  % source: base_profit_prior_weak (4.80), base_profit_prior_strong (4.29),
  % base_spread_weak (0.0167), base_spread_strong (0.667)
  Against the weak incumbent the investor does not trade and the price is
  uninformative. Against the strong incumbent it trades fully and the price is
  informative.

  Now the bottom. Entry goes from a quarter to just over a half, about
  twenty-seven percentage points.
  % source: base_entry_weak (0.250), base_entry_strong (0.523), base_entry_change_pp (27.28)
  High-value ownership goes from one eighth to just under a third.
  % source: base_ownership_weak (0.125), base_ownership_strong (0.324)
  Extra entry follows good news, which is likelier when the challenger is high
  value, so ownership rises too.

  So the challenger's prize shrinks, and entry rises. The price, not the prize,
  brings it in. These are unique equilibrium outcomes at declared parameters;
  the units are arbitrary. \Advance
  \Cue{Now the numbers. \textit{End: advance to the same title with the $r_2$ column.}}
}

\scriptframe{At the benchmark, entry rises from 0.250 to 0.523}{about 1.0 min}{%
  Now add a third column: an incumbent stronger still, r-two, at 3.6.
  % source: base_r_collapse (3.6)
  Profit at the prior is lower still, and the spread is wider still. Trading
  stays fully informative. But entry is back at a quarter, and high-value
  ownership is back at one eighth.
  % source: base_entry_collapse (0.250); ownership 0.125 from tables/table2_equilibrium_controls.tex Panel A

  Why? We saw it on the profit figure. Past a ceiling strength, even the best
  price cannot cover expensive preparation, so only cheap challengers enter,
  informative price or not. Any stronger incumbent at which the floor and
  profitable trading persist, but the best price falls short, gives the same
  fall; 3.6 is simply the declared node.

  So the effect is rise, then fall. It is not monotone in strength. And the
  reversal compares economies in which trading is unique, so no equilibrium
  selection is needed.
  \Cue{Now add a third column. \textit{(Build: $r_2$ column.)}}
}

% ======================================================================
\scriptframe{Freeze the information and deterrence returns}{about 2.3 min}{%
  This is the slide to come back to whenever someone doubts the channel. Does
  the rise really come from information? Orient the table: two strengths, a
  direction and a status.

  The top row is the equilibrium you just saw: price seen, trading re-solved in
  each economy. Entry rises, from a quarter to just over a half.
  % source: base_entry_weak (0.250), base_entry_strong (0.523)

  Now the information controls. First, freeze the investor's orders at full
  size in both economies, so the price carries the same information in both.
  Entry now \emph{falls}, from about 0.56 to about 0.52.
  % source: base_frozen_entry_weak (0.562), base_frozen_entry_strong (0.523)
  That is the textbook deterrence force, back in plain sight. The numbers are a
  numerical diagnostic; the sign is a theorem. Proposition A.3 says that at any
  fixed information experiment, a stronger incumbent cannot raise entry.

  Where does the difference come from? In the strong economy the frozen profile
  is the equilibrium, so the two entries coincide. The whole gap is in the weak
  economy. There, full orders would lose money, so the frozen weak profile is a
  control, not an equilibrium. Its job is to hold information fixed and isolate
  deterrence. If you believed deterrence were the whole story, this is the row
  you would expect to see, and you do see it once the information is held
  fixed. That is why I call the main result a reversal of the received
  comparative static, not an absence of deterrence.

  Second, hide the price from the challenger. Entry is a quarter at both
  strengths.
  % source: base_hidden_entry_weak (0.250), base_hidden_entry_strong (0.250)
  This row is not a second test. It is implied by the high-cost window: without
  the price, expensive preparation never pays, so only the entry floor remains.
  What it shows is that the rise needs the price.

  So hold the price's information fixed and a stronger incumbent weakly lowers
  entry, as the textbook says. The rise comes from what the price reveals.
  \Cue{This is the slide to come back to whenever someone doubts the channel.}
}

% ======================================================================
\scriptframe{At intermediate strength, equilibria coexist}{about 1.6 min}{%
  What happens between the weak and the strong incumbent? The honest answer is:
  partly known. At three strengths, 1.55, 1.60 and 1.65, I can prove that an
  informative equilibrium exists.
  % source: cert_a_r, cert_b_r, cert_c_r
  The proof is computer-assisted: interval arithmetic on exact decimal inputs.

  In each one the investor buys fully after good news and sells partially after
  bad news. The partial sale grows with strength, from a bit under half a unit
  to about nine tenths.
  % source: cert_a_v_interval (q_L about -0.460), cert_c_v_interval (q_L about -0.903), negated
  Certified entry is about 0.55, rising slightly from one strength to the next,
  and above its level at r-one.
  % source: cert_a_entry_interval (0.545), cert_b_entry_interval (0.549), cert_c_entry_interval (0.551); base_entry_strong (0.523)

  Each of these economies also has a no-trade equilibrium, with entry at the
  floor of a quarter.
  % source: base_rho (0.25)
  Why can both exist? Entry is one of the three factors in the investor's
  advantage. When the price is informative, good news draws in the expensive
  challenger, which raises the advantage and supports the trade. When the price
  is flat, only cheap challengers enter, and a single order cannot cover its
  cost. So here the two equilibria coexist.

  Elsewhere the numerical search is not exhaustive, and the full correspondence
  is open. I report every branch I can certify and read no monotone path into
  the gaps. The main theorem does not need this region: it compares economies
  where trading is unique.
  \Cue{What happens between the weak and the strong incumbent?}
}

% ======================================================================
\scriptframe{The entry reversal survives four model changes}{about 1.4 min}{%
  Is the benchmark doing the work? Each row here is its own analytical result,
  with its own declared parameters, so compare within a row, not across rows.

  Atomless costs, a smooth cost distribution instead of two atoms: the numbers
  barely move. The cost atoms are not doing the work.
  % source: cost_mix_laplace_entry_strong (0.523)

  Logistic noise: same sign, smaller size, with entry about 0.3 against the
  strong incumbent.
  % source: logistic_entry_strong (0.302)
  Why does it shrink? With logistic noise the posterior only approaches its
  bound in the tails, so good news is less decisive.

  A small value gap, h equal to two instead of ten: the rise is about as large
  as in the benchmark.
  % source: moderate_h (2), base_h (10), moderate_entry_weak (0.250), moderate_entry_strong (0.527)

  And complementary signals. Here the challenger has its own signal, more
  accurate than the investor's, 0.75 against 0.70. Entry still rises, from 0.85
  to about 0.88.
  % source: signal_buyer_accuracy_value (0.75), signal_trader_accuracy_value (0.70),
  % signal_entry_weak (0.850), signal_entry_strong (0.879)
  The price matters when the challenger's own signal is informative but not
  decisive. With a very accurate own signal, the challenger already enters
  against the weak incumbent, and then entry falls by a few points.
  % source: tables/table_signal_grid.tex (falls about 4 pp when challenger accuracy >= 0.76)

  So the reversal from r-zero to r-one survives all four. I do not claim that
  the fall at r-two survives them.
  \Cue{Is the benchmark doing the work?}
}

% ======================================================================
\scriptframe{Price access raises proceeds and surplus at $r_1$}{about 1.0 min}{%
  One welfare comparison. Hold the incumbent at r-one, and let the challenger
  either see the price or not. Trading is re-solved in both, and the investor
  trades fully in both, so trading costs coincide.

  Seeing the price raises expected target proceeds from about 0.61 to about
  0.87, a gain of about a quarter.
  % source: base_revenue_hidden (0.615), base_revenue_feedback (0.872), base_revenue_gain (0.258)
  Acquisition surplus net of preparation costs rises by about 0.08.
  % source: base_net_surplus_gain (0.0802)
  The surplus gain is not an accident of accounting: each extra entry happens
  only when its expected profit covers its cost, and the gain also includes
  sales that would otherwise fail the reserve.

  At r-zero, prices are uninformative, so there is no gain. This is a
  comparison of economies at r-one. It does not rank incumbent strengths, and it
  is not a recommendation on process design.
  \Cue{One welfare comparison.}
}

% ======================================================================
\scriptframe{What an empirical test would have to measure}{about 1.1 min}{%
  Can you test this? The paper gives a design, not a result. There are four
  requirements.

  The outcome has to be the decision to prepare: the start of substantive
  diligence, or a proposal that needs costly preparation. Counting public offers
  is not enough. Strength has to be measured from information that was public
  before that decision. The price has to come before preparation; a stock that
  merely trades during confidential negotiations does not qualify.

  And the hard one: target returns can reflect anticipated bidder arrival. So
  returns followed by entry do not, by themselves, show that anyone learned from
  prices. The online appendix proposes a pilot that rebuilds these chronologies
  from disclosure records, separating when an event happened from when it
  became public. That is why timing has to come from the records, not from
  announcement dates. There is no sample and no estimate yet.
  \Cue{Can you test this?}
}

% ======================================================================
\begin{samepage}
\scriptframe{Conclusion}{about 0.9 min}{%
  Let me close with the answer. A stronger incumbent can bring the challenger
  in. At the benchmark, entry roughly doubles, from a quarter to just over a
  half, while the challenger's expected profit falls.
  % source: base_entry_weak (0.250), base_entry_strong (0.523)

  The reason is one auction with two claims. Competition makes target proceeds
  more sensitive to who the challenger is. Informed trading puts that into the
  price, and good news draws in the expensive challenger. Freeze that
  information, and deterrence returns.

  The next theorem belongs to the seller: which sale terms to choose, when terms
  move both what the winner pays and what the price reveals.

  So when the target trades, a stronger rival is not a pure deterrent. Who
  competes depends on what the price reveals before anyone prepares.
  \Cue{Let me close with the answer.}
}
\end{samepage}

\Needspace{14\baselineskip}
\section*{Live cue layer}
{\raggedright\small
\begin{enumerate}\setlength{\itemsep}{0.25em}
  \item \textbf{Opening (title)} --- Thank you for having me.
  \item \textbf{Does a stronger incumbent keep the challenger out?} --- Here is the situation the paper is about.
  \item \textbf{This paper} --- Here is what I do and what I find.
  \item \textbf{What is new relative to learning from prices} --- Three literatures meet here.
  \item \textbf{Model: who moves, who knows what, and when} --- Here is the model, in one timeline.
  \item \textbf{One auction, two claims, opposite responses} --- This table is the whole paper in one place.
  \item \textbf{A stronger incumbent: spread up, profit down} --- Here are the two claims at benchmark scale.
  \item \textbf{Trading pays only against a strong incumbent} --- Now the investor.
  \item \textbf{Only good news makes expensive preparation pay} --- Now the challenger's side: what does a moving price tell it?
  \item \textbf{Proposition 2: a stronger incumbent can raise entry} --- Here is the main theorem, in words first.
  \item \textbf{At the benchmark, entry rises from 0.250 to 0.523} --- Now the numbers. \textit{[advance at the end]}
  \item \textbf{At the benchmark, entry rises from 0.250 to 0.523} ($r_2$ column) --- Now add a third column.
  \item \textbf{Freeze the information and deterrence returns} --- This is the slide to come back to whenever someone doubts the channel.
  \item \textbf{At intermediate strength, equilibria coexist} --- What happens between the weak and the strong incumbent?
  \item \textbf{The entry reversal survives four model changes} --- Is the benchmark doing the work?
  \item \textbf{Price access raises proceeds and surplus at $r_1$} --- One welfare comparison.
  \item \textbf{What an empirical test would have to measure} --- Can you test this?
  \item \textbf{Conclusion} --- Let me close with the answer.
\end{enumerate}\par}

% ######################################################################
\ScriptAppendix

\scriptframe{Which deals have a decision interval}{if asked (A1)}{%
  If asked whether this is realistic, jump to A1. Whether a deal fits is a
  question about its chronology. A disclosed approach, an announced strategic
  review or an open contest can create an interval; none does so automatically.
  A wholly confidential process does not, because no challenger could watch the
  price while it mattered. Strength is the public distribution of the
  incumbent's value, not an announced bid. The paper settles this for no
  specific transaction.
}

\scriptframe{Evidence: a design, not a result}{if asked (A2)}{%
  If asked for evidence: the paper has a design, not a result. The pilot
  rebuilds from disclosure records when approaches, contacts, diligence and
  proposals happened, and separately when each became public. The outcome is
  the decision to prepare, and prices must come before it. The hard part is that
  later returns can reflect anticipated arrival. There is no sample, effect or
  identification strategy yet.
}

\scriptframe{Related work in more detail}{if asked (A3)}{%
  If asked how this differs from feedback models: in Dow, Goldstein and Guembel
  the price guides the firm's own investment. Here the sale rule splits one
  surplus into a traded claim and the challenger's claim, and competition moves
  them in opposite directions. The informed party is not a bidder, and it trades
  before entry. Only after checking both papers, if Kyle and Vila or Schwert come
  up: there the informed trader is the raider, and run-ups follow an anticipated
  bid. Here the price moves before the challenger decides.
}

\scriptframe{What the model leaves out, and why}{if asked (A4)}{%
  If asked about toeholds or an incumbent moving the price: neither bidder
  trades in the model, and the investor knows theta, not the incumbent's value.
  Toeholds, free riding, jump bids as signals, and incumbent or target
  manipulation are outside the model. The paper does not solve toeholds, so I
  do not guess the direction of their effect. Inside the model, the investor
  cannot profit by faking good news.
}

\scriptframe{Complementary, not superior, information}{if asked (A5)}{%
  If asked why an investor would know more than the challenger: it need not
  know more; it knows different things. The challenger knows its integration
  plans; a specialist investor may know the target's customers, technology and
  demand. Preparation is what reveals the exact value. The benchmark's perfect
  investor only makes the mechanism visible. The reversal holds when the
  challenger's own signal is more accurate, 0.75 against 0.70.
  % source: signal_buyer_accuracy_value (0.75), signal_trader_accuracy_value (0.70)
}

\scriptframe{The price helps a challenger with its own signal}{if asked (A6)}{%
  If pressed on the signal economy: the challenger combines its own signal with
  the price. In the declared example entry rises from 0.85 to about 0.88, and
  that is analytical.
  % source: signal_entry_weak (0.850), signal_entry_strong (0.879)
  Across the twenty-five-cell grid, entry rises in the six cells that meet the
  condition, is unchanged in nine, and falls by about four points when the
  challenger's signal is very accurate. Those last two are numerical
  diagnostics. The price matters when the own signal is informative but not
  decisive.
  % source: tables/table_signal_grid.tex
}

\scriptframe{Equilibrium and the two comparisons}{if asked (A7)}{%
  If asked about the equilibrium concept: order distributions after each value,
  with mixing allowed; a rational price that anticipates entry; Bayesian beliefs
  from the price; optimal entry; truthful bids. Keep two comparisons apart. A
  deviation is evaluated against fixed price and entry schedules. A comparison
  across economies re-solves everything. Uniqueness refers to trading and
  on-path entry under truthful bidding.
}

\scriptframe{Acquisition payoffs in closed form}{if asked (A8)}{%
  If asked whether this is a uniform-distribution artifact: no. The closed forms
  are the uniform case, but Proposition 1 holds for any incumbent distribution
  whose top lies between ell and h. The spread is the expected amount by which
  the incumbent's value exceeds ell, which rises under any first-order
  strengthening, while the challenger's profit falls.
}

\scriptframe{Why trading is unique: a global bound}{if asked (A10)}{%
  If asked whether I assumed the informative profile: no. Against the weak
  incumbent the gross advantage per unit is capped by the spread, which is below
  k, so every nonzero order loses. Against the strong incumbent the marginal
  profit of every unit exceeds k under every candidate price and entry schedule,
  mixed orders included. It is a global bound, not a first-order condition. The
  same inequalities make a wrong-signed order lose, so faking good news does not
  pay.
}

\scriptframe{Orders, units and the trading cost}{if asked (A11)}{%
  If asked why a corner, unlike Kyle: the per-unit advantage never falls below
  rho times m times the spread, because beliefs stay in the band. And one more
  unit erodes it by at most a fraction one over b. So marginal profit stays above
  k across the whole interval. The order bound is a normalized trading unit, not
  ownership, and k is an execution friction separate from price impact. If k
  were near zero, trading would pay against the weak incumbent too, and the
  theorem is silent.
}

\scriptframe{The challenger can read the market's belief off the price}{if asked (A13)}{%
  If asked whether the challenger can really read the belief off the price: yes.
  The price is proceeds without a challenger plus entry probability times a term
  that rises with the belief. Because entry is at least rho and the spread is
  positive, the price strictly increases in the belief, even with atoms and
  entry jumps. So the price is a sufficient statistic for the market's belief,
  and seeing the order flow would add nothing.
}

\scriptframe{Why Laplace noise, and what logistic noise changes}{if asked (A14)}{%
  If asked whether Laplace noise does the work: Laplace and logistic noise
  satisfy the same derivative bound, so the global trading bounds hold for both.
  The difference is where the mass sits. Under Laplace the posterior reaches its
  bound at finite flow; under logistic, only in the limit. So good news is less
  decisive, and entry against the strong incumbent is about 0.3 instead of just
  over a half. Same sign, smaller size, and not a Blackwell ranking.
  % source: logistic_entry_strong (0.302), base_entry_strong (0.523)
}

\scriptframe{Why the model needs a low-cost floor}{if asked (A15)}{%
  If asked whether the cheap-type floor does all the work: the floor does three jobs in the proofs. Entry is at least rho after every price, so the price reveals the market's posterior, and the investor's edge never falls below rho m Delta-T, which is what beats the trading cost. And cheap preparation does not respond to the price, so information can only add entry. If cheap preparation paid only at good prices, bad prices would deter it and flows could pool at one price. If it did not pay even at the prior, no preparation and no trade would be an equilibrium. But the floor alone gives entry of a quarter at both
  strengths when the price is hidden. It is necessary, and it does not produce
  the rise by itself.
  % source: base_hidden_entry_weak, base_hidden_entry_strong (0.250)
}

\scriptframe{Proof logic in four steps}{if asked (A16)}{%
  If asked what does the work in the proof: step two, the global bound that
  forces full orders against the strong incumbent. Step one rules out trading
  against the weak incumbent. Step three uses the high-cost window, so good news
  crosses the entry threshold, and more often when the challenger is high value.
  Step four is Blackwell: the strong price experiment strictly dominates the
  uninformative one.
}

\scriptframe{Conditions are jointly satisfiable for every $h>\ell$}{if asked (A17)}{%
  If asked whether this is a knife edge: the smallest benchmark margin is small,
  about a quarter of a hundredth against a trading cost of two hundredths.
  % source: base_minimum_theorem_margin (0.00241), base_k (0.02)
  But nonemptiness for every h above ell does not come from slack at the
  benchmark. It comes from a construction: put the two strengths close to ell,
  then place k and the costs strictly inside their windows. The conditions are
  strict inequalities, so the set is open. Part three is not claimed for every h.
}

\scriptframe{Two forces from one payment rule}{if asked (A18)}{%
  If asked which force wins: put them side by side. Deterrence lowers the
  challenger's profit at every belief; it needs nothing, and at fixed
  information it weakly lowers entry. Information widens the spread; it needs
  the low-cost floor and a price seen before entry, and it raises entry once
  trading pays and good news clears the expensive cost. The three named
  conditions decide which one wins.
}

\scriptframe{Benchmark scale: declared, not calibrated}{if asked (A19)}{%
  If asked whether the magnitudes are calibrated: no. The units are arbitrary,
  and the benchmark is declared to show that the conditions hold together and to
  scale the effect. The signs are the theorem. A moderate economy, with h equal
  to two and much smaller costs, gives a rise of similar size, from a quarter to
  about 0.53. So the tenfold value gap is not the mechanism.
  % source: moderate_h (2), moderate_entry_weak (0.250), moderate_entry_strong (0.527)
}

\scriptframe{How entry and ownership are computed}{if asked (A20)}{%
  If asked why ownership rises, or why entry falls at r-two: the expensive
  challenger enters when order flow crosses a threshold. High-value challengers
  cross it more often, so the extra entry tilts toward them and ownership rises.
  Above the ceiling strength, about 3.6, the entry threshold exceeds the best
  belief a price can induce, so no order flow brings the expensive challenger
  in.
  % source: base_r_high_cost_ceiling (3.59)
}

\scriptframe{Which equilibrium? What the proof certifies}{if asked (A21)}{%
  If asked which equilibrium I select: none. The reversal compares economies
  where trading and on-path entry are unique, so there is nothing to select;
  multiplicity appears only in between. At the three certified strengths,
  interval arithmetic brackets the low type's optimal order and verifies that
  the high type's full purchase dominates. Not proved: uniqueness of the
  informative profile, absence of mixed equilibria, or a branch between nodes.
}

\scriptframe{Information controls in detail}{if asked (A22)}{%
  If asked whether the frozen control is fair: it is a control, not an
  equilibrium at r-zero, and that is its purpose. It holds the information
  experiment fixed so that only deterrence moves. Its numbers are a numerical
  diagnostic; its sign is Proposition A.3. The hidden-price row is implied by
  the high-cost window, so it is not a test of deterrence. It shows that without
  the price the rise disappears.
}

\scriptframe{Price level or information?}{if asked (A23)}{%
  If asked whether this is just a higher price level: no. Add a dividend equal
  to the revenue gain, about a quarter, to the traded claim in the price-hidden
  economy. Mean prices then match, but entry does not, so what matters is the
  information in the price, not its level. That is a diagnostic; the revenue
  gain itself is analytical. The dividend is not a sale mechanism, and I do not
  rank strengths or mechanisms.
  % source: base_matched_dividend (0.258), base_revenue_gain (0.258)
}

\scriptframe{Trading and entry across strengths: entry}{if asked (A24)}{%
  If asked what the whole correspondence looks like: no trade is unique for weak
  enough incumbents and remains an equilibrium up to about 1.75. Full orders are
  unique above about 2.84. Expensive entry is impossible above about 3.6. In
  between, the black points are certified and the curves are numerical
  diagnostics, broken where unresolved. The search is not exhaustive, and the
  correspondence is open.
  % source: base_r_no_trade_exact (1.75), base_r_full_unique_sufficient (2.84), base_r_high_cost_ceiling (3.59)
}

\scriptframe{Robustness in full}{if asked (A26)}{%
  If asked whether the rows share parameters: no. Each row has its own declared
  parameters, so the table is not one joint calibration. Every row gives the
  rise from r-zero to r-one analytically, and the smallest theorem margins are
  listed underneath. The fall at r-two is shown for the benchmark only.
}

\scriptframe{A higher reserve can raise proceeds; optimal terms are open}{if asked (A27)}{%
  If asked what the seller should do: I do not have the optimal terms. At the
  listed nodes, raising the reserve from one half to 1.1 raises expected
  proceeds against both incumbents.
  % source: base_p (0.5), value_reserve_high (1.1), value_revenue_* (0.393 to 0.432; 0.872 to 1.01)
  The value 1.1 sits just above ell, so it excludes the low-value challenger,
  and preparation, sale and two admissible bidders come apart. That is a
  feasible improvement, not an optimal reserve. The seller's problem is the next
  theorem.
}

\scriptframe{Same orders, different prices}{if asked (A28)}{%
  If asked why the seller's problem is hard: the same orders can support
  different prices. At a high reserve, full orders support price pools below
  different cutoffs, with different pooled beliefs, entry and proceeds. So a
  continuation has to specify the price rule, not just the orders. This does not
  arise on the benchmark support.
}

\Needspace{12\baselineskip}
\section*{Standing Q\&A stance}
Hear the whole question and count to three. Give the economic answer first,
then the backup slide only if it makes the answer more precise. Explain rather
than capitulate; concede a genuine scope condition directly, bad news first and
then the good news. Keep the status words the slides use: analytical,
computer-assisted, numerical diagnostic, open.

\textit{Why a second-price auction? What about negotiation?} (frame 5, then A9;
kept as prose here because A9 carries no script block): under Nash bargaining
with seller weight eta, challenger profit still falls weakly with strength for
every eta, but the spread effect rises weakly below one half and falls weakly
above. So the opposition itself depends on the payment rule. No entry reversal
is claimed under bargaining, and first-price equilibria are not solved.

Answers that live on main frames, with no backup: \textit{Why can't the challenger
bid without preparing?} (frames 1 and 4: there is no executable bid without
preparation, so a challenger that declines stays out). \textit{What if the
challenger saw order flow?} (A13: the price already carries the market's belief).
\textit{Does the fall at $r_2$ survive the extensions?} (frame 13: shown at the
benchmark and nearby only). \textit{Why must the price come before preparation?}
(frame 4: otherwise there is nothing to learn; frame 11: hiding the price leaves
only the floor).



\end{document}
